\section{Limitations and conclusion}
\label{sec:conclusion}

\paragraph{Limitations.}
A digital zoom is an exact camera change, but real zoom lenses add distortion and change focus; our
optical-zoom control isolates detail loss from the camera change but is still a resampled image. For a
model without camera input, a crop also removes context, so $\alpha=0$ is a target, not a guarantee, on
every image (\cref{sec:audit}). The label-free score predicts zoom-induced error, not accuracy in general,
and fails on DIODE-outdoor. EFT+B trades accuracy on unzoomed and unseen images for equivariance in our
setting; it is a demonstration that the law can train a model, not a drop-in replacement. Stale
intrinsics in camera-conditioned models are a known deployment problem~\cite{kim2026mismatch}; we only
place them on the same scale as the other groups.

\paragraph{Conclusion.}
A crop is the simplest change of camera there is, and its effect is known exactly. Metric depth models
that do not receive the camera largely ignore it: they read depth from apparent size, which a zoom cannot
change, and the zoom--squash ambiguity says why. The two exponents make this visible without labels,
place every model on a common scale, and point to what a correction has to avoid: a model can learn to
detect a zoom from blur instead of geometry. We hope the audit becomes a routine check for any model that
claims to output metres.
